\documentclass[11pt,a4paper]{article}
\usepackage[utf8]{inputenc}
\usepackage[T1]{fontenc}
\usepackage[spanish,es-noshorthands]{babel}
\usepackage[margin=2.5cm]{geometry}
\usepackage{amsmath,amssymb,mathtools}
\usepackage{enumitem}
\usepackage{booktabs}
\usepackage{tikz}
\usepackage{pgfplots}
\pgfplotsset{compat=1.18}

\newcounter{ejnum}
\newenvironment{ejercicio}{\refstepcounter{ejnum}\par\medskip\noindent\textbf{Ejercicio \theejnum.}~}{\par}
\newenvironment{solucion}{\par\noindent\textit{Solución.}~}{\par\vspace{1em}}

\title{Práctica 3: Bondad de ajuste, independencia y homogeneidad ($\chi^2$) \\ \large Unidad 7: Prueba de hipótesis}
\author{Probabilidad y Estadística (5765) \\ Universidad Nacional del Sur}
\date{}

\begin{document}
\maketitle

\begin{ejercicio}
Una empresa de encuestas afirma que las preferencias de los consumidores por cuatro marcas de gaseosa (A, B, C, D) se reparten en las proporciones $0.40,\ 0.30,\ 0.20,\ 0.10$ respectivamente. Se encuesta a $200$ consumidores y se obtienen las siguientes frecuencias observadas: A: $70$, B: $65$, C: $40$, D: $25$. Con $\alpha=0.05$, pruebe si las preferencias observadas son compatibles con las proporciones declaradas.
\end{ejercicio}

\begin{solucion}
\textbf{Hipótesis:} $H_0: p_A=0.40,\ p_B=0.30,\ p_C=0.20,\ p_D=0.10$ vs. $H_1:$ al menos una proporción difiere.

\textbf{Frecuencias esperadas} ($E_i=200\,p_i$): $E_A=80$, $E_B=60$, $E_C=40$, $E_D=20$. Todas $\ge5$: se cumple la condición de aplicación.

\textbf{Grados de libertad:} $k-1-m=4-1-0=3$ (no se estimó ningún parámetro).

\textbf{Región crítica:} $\chi^2>\chi^2_{3,\,0.05}=7.815$.

\textbf{Cálculo:}
\[
\chi^2=\frac{(70-80)^2}{80}+\frac{(65-60)^2}{60}+\frac{(40-40)^2}{40}+\frac{(25-20)^2}{20}=\frac{100}{80}+\frac{25}{60}+0+\frac{25}{20}=1.25+0.417+0+1.25=2.92.
\]

\textbf{Decisión:} $2.92<7.815\Rightarrow$ \textbf{no se rechaza} $H_0$.

\textbf{Conclusión:} al $5\%$ de significación, las preferencias observadas son compatibles con las proporciones declaradas por la empresa.
\end{solucion}

\begin{ejercicio}
Un genetista cruza dos variedades de plantas y espera, según la teoría mendeliana, que la descendencia se reparta en la proporción $9:3:3:1$ entre cuatro fenotipos. Sobre $160$ plantas observadas se registraron las frecuencias: $92,\ 34,\ 24,\ 10$. Con $\alpha=0.05$, ¿son los datos compatibles con la proporción teórica $9:3:3:1$?
\end{ejercicio}

\begin{solucion}
\textbf{Hipótesis:} $H_0: (p_1,p_2,p_3,p_4)=\left(\tfrac{9}{16},\tfrac{3}{16},\tfrac{3}{16},\tfrac{1}{16}\right)$ vs. $H_1:$ no se cumple esa proporción.

\textbf{Frecuencias esperadas:} $E_1=160\times\tfrac{9}{16}=90$, $E_2=160\times\tfrac{3}{16}=30$, $E_3=30$, $E_4=160\times\tfrac{1}{16}=10$. Todas $\ge5$.

\textbf{Grados de libertad:} $4-1-0=3$.

\textbf{Región crítica:} $\chi^2>\chi^2_{3,\,0.05}=7.815$.

\textbf{Cálculo:}
\[
\chi^2=\frac{(92-90)^2}{90}+\frac{(34-30)^2}{30}+\frac{(24-30)^2}{30}+\frac{(10-10)^2}{10}=\frac{4}{90}+\frac{16}{30}+\frac{36}{30}+0=0.044+0.533+1.2+0=1.78.
\]

\textbf{Decisión:} $1.78<7.815\Rightarrow$ \textbf{no se rechaza} $H_0$.

\textbf{Conclusión:} al $5\%$ de significación, los datos observados son compatibles con la proporción mendeliana $9:3:3:1$.
\end{solucion}

\begin{ejercicio}
El número de llamadas que recibe un call center por minuto, registrado durante $150$ minutos, se resume en la siguiente tabla. Se desea comprobar si el número de llamadas por minuto sigue una distribución de Poisson con media estimada a partir de los datos, $\hat\lambda=2.0$.
\begin{center}
\begin{tabular}{@{}lcccccc@{}}
\toprule
N° de llamadas & 0 & 1 & 2 & 3 & 4 & $\ge5$ \\
\midrule
$O_i$ (minutos) & 22 & 40 & 38 & 28 & 15 & 7 \\
\bottomrule
\end{tabular}
\end{center}
Las probabilidades de Poisson($\lambda=2.0$) para cada categoría son $p_0=0.135$, $p_1=0.271$, $p_2=0.271$, $p_3=0.180$, $p_4=0.090$, $p_{\ge5}=0.053$. Con $\alpha=0.05$, pruebe la bondad del ajuste.
\end{ejercicio}

\begin{solucion}
\textbf{Hipótesis:} $H_0:$ el número de llamadas por minuto sigue una distribución Poisson($\lambda=2.0$) vs. $H_1:$ no sigue esa distribución.

\textbf{Frecuencias esperadas} ($E_i=150\,p_i$):
\[
E_0=20.25,\quad E_1=40.65,\quad E_2=40.65,\quad E_3=27.0,\quad E_4=13.5,\quad E_{\ge5}=7.95.
\]
Todas $\ge5$: se cumple la condición.

\textbf{Grados de libertad:} $k-1-m=6-1-1=4$ (se estimó $\lambda$ a partir de los datos, $m=1$).

\textbf{Región crítica:} $\chi^2>\chi^2_{4,\,0.05}=9.488$.

\textbf{Cálculo:}
\[
\chi^2=\frac{(22-20.25)^2}{20.25}+\frac{(40-40.65)^2}{40.65}+\frac{(38-40.65)^2}{40.65}+\frac{(28-27.0)^2}{27.0}+\frac{(15-13.5)^2}{13.5}+\frac{(7-7.95)^2}{7.95}
\]
\[
\approx0.151+0.010+0.173+0.037+0.167+0.114=0.65.
\]

\textbf{Decisión:} $0.65<9.488\Rightarrow$ \textbf{no se rechaza} $H_0$.

\textbf{Conclusión:} al $5\%$ de significación, el número de llamadas por minuto es compatible con un modelo de Poisson de media $2.0$.
\end{solucion}

\begin{ejercicio}
En un hospital se estudia si el tipo de complicación postoperatoria (leve, moderada, grave) está relacionado con el tipo de cirugía realizada (programada o de urgencia), sobre una muestra de $250$ pacientes:
\begin{center}
\begin{tabular}{@{}lccc|c@{}}
\toprule
 & Leve & Moderada & Grave & Total \\
\midrule
Programada & 78 & 42 & 10 & 130 \\
Urgencia & 52 & 48 & 20 & 120 \\
\midrule
Total & 130 & 90 & 30 & 250 \\
\bottomrule
\end{tabular}
\end{center}
Con $\alpha=0.05$, ¿hay evidencia de que el tipo de complicación depende del tipo de cirugía?
\end{ejercicio}

\begin{solucion}
\textbf{Hipótesis:} $H_0:$ tipo de complicación y tipo de cirugía son independientes vs. $H_1:$ no son independientes.

\textbf{Frecuencias esperadas} ($E_{ij}=O_{i\bullet}O_{\bullet j}/250$):
\[
E_{11}=\frac{130\times130}{250}=67.6,\ \ E_{12}=\frac{130\times90}{250}=46.8,\ \ E_{13}=\frac{130\times30}{250}=15.6,
\]
\[
E_{21}=\frac{120\times130}{250}=62.4,\ \ E_{22}=\frac{120\times90}{250}=43.2,\ \ E_{23}=\frac{120\times30}{250}=14.4.
\]
Todas $\ge5$: se cumple la condición de aplicación.

\textbf{Grados de libertad:} $(r-1)(c-1)=(2-1)(3-1)=2$.

\textbf{Región crítica:} $\chi^2>\chi^2_{2,\,0.05}=5.991$.

\textbf{Cálculo:}
\[
\chi^2=\frac{(78-67.6)^2}{67.6}+\frac{(42-46.8)^2}{46.8}+\frac{(10-15.6)^2}{15.6}+\frac{(52-62.4)^2}{62.4}+\frac{(48-43.2)^2}{43.2}+\frac{(20-14.4)^2}{14.4}
\]
\[
=\frac{108.16}{67.6}+\frac{23.04}{46.8}+\frac{31.36}{15.6}+\frac{108.16}{62.4}+\frac{23.04}{43.2}+\frac{31.36}{14.4}
\]
\[
=1.600+0.492+2.010+1.733+0.533+2.178=8.55.
\]

\textbf{Decisión:} $8.55>5.991\Rightarrow$ se \textbf{rechaza} $H_0$.

\textbf{Conclusión:} al $5\%$ de significación, hay evidencia de que el tipo de complicación postoperatoria \textbf{no} es independiente del tipo de cirugía; las cirugías de urgencia muestran proporcionalmente más complicaciones graves de lo esperado bajo independencia.
\end{solucion}

\begin{ejercicio}
Se investiga si existe relación entre el nivel de estudios de los padres (secundario, terciario/universitario) y si sus hijos practican algún deporte federado (sí/no), sobre una muestra de $180$ familias:
\begin{center}
\begin{tabular}{@{}lcc|c@{}}
\toprule
 & Practica deporte & No practica & Total \\
\midrule
Secundario & 38 & 62 & 100 \\
Terciario/Univ. & 42 & 38 & 80 \\
\midrule
Total & 80 & 100 & 180 \\
\bottomrule
\end{tabular}
\end{center}
Con $\alpha=0.05$, pruebe si la práctica deportiva de los hijos es independiente del nivel de estudios de los padres.
\end{ejercicio}

\begin{solucion}
\textbf{Hipótesis:} $H_0:$ las variables son independientes vs. $H_1:$ no son independientes.

\textbf{Frecuencias esperadas:}
\[
E_{11}=\frac{100\times80}{180}=44.4,\quad E_{12}=\frac{100\times100}{180}=55.6,\quad E_{21}=\frac{80\times80}{180}=35.6,\quad E_{22}=\frac{80\times100}{180}=44.4.
\]

\textbf{Grados de libertad:} $(2-1)(2-1)=1$.

\textbf{Región crítica:} $\chi^2>\chi^2_{1,\,0.05}=3.841$.

\textbf{Cálculo:}
\[
\chi^2=\frac{(38-44.4)^2}{44.4}+\frac{(62-55.6)^2}{55.6}+\frac{(42-35.6)^2}{35.6}+\frac{(38-44.4)^2}{44.4}
\]
\[
=\frac{40.96}{44.4}+\frac{40.96}{55.6}+\frac{40.96}{35.6}+\frac{40.96}{44.4}=0.922+0.737+1.151+0.922=3.73.
\]

\textbf{Decisión:} $3.73<3.841\Rightarrow$ \textbf{no se rechaza} $H_0$ (aunque el valor está muy cerca del límite crítico).

\textbf{Conclusión:} al $5\%$ de significación, no hay evidencia suficiente para afirmar que la práctica deportiva de los hijos esté relacionada con el nivel de estudios de los padres; el resultado, muy próximo al valor crítico, sugiere ampliar la muestra para obtener una conclusión más robusta.
\end{solucion}

\begin{ejercicio}
Se comparan tres turnos de una fábrica (mañana, tarde, noche) respecto de la proporción de piezas que resultan aptas o no aptas en el control de calidad. Se extrae una muestra fija de $120$ piezas de cada turno:
\begin{center}
\begin{tabular}{@{}lccc|c@{}}
\toprule
 & Mañana & Tarde & Noche & Total \\
\midrule
Aptas & 108 & 102 & 90 & 300 \\
No aptas & 12 & 18 & 30 & 60 \\
\midrule
Total & 120 & 120 & 120 & 360 \\
\bottomrule
\end{tabular}
\end{center}
Con $\alpha=0.01$, pruebe si los tres turnos son homogéneos respecto de la calidad de las piezas producidas.
\end{ejercicio}

\begin{solucion}
\textbf{Hipótesis:} $H_0:$ los tres turnos tienen la misma distribución de calidad (aptas/no aptas) vs. $H_1:$ al menos un turno difiere.

\textbf{Frecuencias esperadas} ($E_{ij}=O_{i\bullet}O_{\bullet j}/360$): fila ``Aptas'': $E=\frac{300\times120}{360}=100$ para cada turno; fila ``No aptas'': $E=\frac{60\times120}{360}=20$ para cada turno.

\textbf{Grados de libertad:} $(r-1)(c-1)=(2-1)(3-1)=2$.

\textbf{Región crítica:} $\chi^2>\chi^2_{2,\,0.01}=9.210$.

\textbf{Cálculo:}
\[
\chi^2=\frac{(108-100)^2}{100}+\frac{(102-100)^2}{100}+\frac{(90-100)^2}{100}+\frac{(12-20)^2}{20}+\frac{(18-20)^2}{20}+\frac{(30-20)^2}{20}
\]
\[
=\frac{64}{100}+\frac{4}{100}+\frac{100}{100}+\frac{64}{20}+\frac{4}{20}+\frac{100}{20}=0.64+0.04+1.00+3.20+0.20+5.00=10.08.
\]

\textbf{Decisión:} $10.08>9.210\Rightarrow$ se \textbf{rechaza} $H_0$.

\textbf{Conclusión:} al $1\%$ de significación, hay evidencia de que los tres turnos \textbf{no} son homogéneos respecto de la calidad de las piezas producidas; el turno noche presenta una proporción de piezas no aptas notablemente mayor a la esperada bajo homogeneidad.
\end{solucion}

\begin{ejercicio}
Un centro de salud compara la efectividad de tres tratamientos distintos contra una misma afección, asignando (por diseño) $80$ pacientes a cada tratamiento y registrando si hubo mejoría o no al cabo de un mes:
\begin{center}
\begin{tabular}{@{}lccc|c@{}}
\toprule
 & Tratamiento 1 & Tratamiento 2 & Tratamiento 3 & Total \\
\midrule
Mejoría & 52 & 61 & 45 & 158 \\
Sin mejoría & 28 & 19 & 35 & 82 \\
\midrule
Total & 80 & 80 & 80 & 240 \\
\bottomrule
\end{tabular}
\end{center}
\begin{enumerate}[label=(\alph*)]
\item Explique por qué esta situación corresponde a una prueba de \textbf{homogeneidad} y no de independencia.
\item Realice la prueba con $\alpha=0.05$ y concluya.
\end{enumerate}
\end{ejercicio}

\begin{solucion}
\textbf{(a)} Corresponde a una prueba de homogeneidad porque los tamaños de cada grupo ($80$ pacientes por tratamiento) fueron \textbf{fijados de antemano} por el diseño del estudio (se decidió cuántos pacientes recibiría cada tratamiento), en lugar de tomar una única muestra aleatoria de pacientes y clasificarla luego según dos variables. Se comparan tres poblaciones (una por tratamiento) respecto de la misma variable categórica (mejoría/sin mejoría).

\textbf{(b) Hipótesis:} $H_0:$ los tres tratamientos tienen igual distribución de mejoría vs. $H_1:$ al menos uno difiere.

\textbf{Frecuencias esperadas:} fila ``Mejoría'': $E=\frac{158\times80}{240}=52.67$ para cada tratamiento; fila ``Sin mejoría'': $E=\frac{82\times80}{240}=27.33$ para cada tratamiento.

\textbf{Grados de libertad:} $(2-1)(3-1)=2$. \textbf{Región crítica:} $\chi^2>\chi^2_{2,\,0.05}=5.991$.

\textbf{Cálculo:}
\[
\chi^2=\frac{(52-52.67)^2}{52.67}+\frac{(61-52.67)^2}{52.67}+\frac{(45-52.67)^2}{52.67}+\frac{(28-27.33)^2}{27.33}+\frac{(19-27.33)^2}{27.33}+\frac{(35-27.33)^2}{27.33}
\]
\[
\approx0.009+1.318+1.116+0.016+2.541+2.151=7.15.
\]

\textbf{Decisión:} $7.15>5.991\Rightarrow$ se \textbf{rechaza} $H_0$.

\textbf{Conclusión:} al $5\%$ de significación, hay evidencia de que los tres tratamientos no son homogéneos en cuanto a la proporción de pacientes que mejoran; el tratamiento 2 muestra una tasa de mejoría notablemente superior a la esperada bajo homogeneidad.
\end{solucion}

\end{document}
